% DISCUSSION AND LIMITATIONS. Source: harness/results/summary/metrics_summary.csv.
\section{Discussion and Limitations}\label{sec:discussion}

Turpin et al.\ reported accuracy drops of up to 36\% under
their biasing features in 2023 models~\cite{turpin2023}. In
our adaptation, which uses only the Suggested Answer cue,
accuracy fell by 1.9 percentage points for GPT-5.6 Sol and
0.4 for Claude Opus 5.5, within sampling error.
Differences in cue types, samples and models rule out a
direct quantitative comparison.

The directional baseline changes how the switches should be read.
Comparing GPT's nine cue-following switches with all seven uncued
switches would suggest little cue effect, but only two uncued
switches landed on the cued option, so GPT's susceptibility,
though rare, is unlikely to be chance. Claude showed no
susceptibility beyond ordinary variability. One repeat per
question still cannot attribute any single switch to the cue.

Where GPT did follow the cue, its explanations never mentioned
it, consistent with the pattern Turpin et al.\ reported and with
the low acknowledgement rates Chen et al.\ found for reasoning
models~\cite{chen2025}. Pooling across models would hide this:
the one explicit acknowledgement came from Claude's single case.

Several limits apply. Only ten cases were available for coding,
nine from one model. The prompt asks for a 2--4 sentence
explanation of ``the key evidence'' and discourages revealing
chain of thought, which may itself reduce mentions of the
suggestion; we measure acknowledgement in that format, not in
free reasoning. We evaluate visible explanations rather than
hidden reasoning, so a \emph{none} label means the suggestion was
not acknowledged in the text, not that the model did or did not
use it. The study covers two models, one cue type and one
reasoning-effort level. Finally, the questions come from public
data~\cite{turpin2023code} and may appear in training data: in
one causal-judgement item, Claude's reasoning summary identified
the scenario as a well-known one. The high Original accuracies
(90.0\% and 96.7\%) may partly reflect such familiarity.
